\uuid{6Jyo}
\exo7id{174}
\titre{exo7 174}
\auteur{cousquer}
\organisation{exo7}
\datecreate{2003-10-01}
\isIndication{false}
\isCorrection{false}
\chapitre{Logique, ensemble, raisonnement}
\sousChapitre{Récurrence}
\cpge{mpsi-mp2i.raisonnement-vocabulaire-ensembliste.a, pcsi.raisonnement-vocabulaire-ensembliste.a, ptsi.raisonnement-vocabulaire-ensembliste.a, mpsi-mp2i.calcul-algebrique-trigonometrie.c, pcsi.calcul-algebrique-trigonometrie.c, ptsi.calcul-algebrique-trigonometrie.c, mpsi-mp2i.calcul-algebrique-trigonometrie.a, pcsi.calcul-algebrique-trigonometrie.a, ptsi.calcul-algebrique-trigonometrie.a, bcpst1.logique-ensembles-raisonnement.c, bcpst1.methodes-calcul}
\module{Algèbre}
\niveau{L1}
\difficulte{}

\contenu{
\begin{enumerate}
    \item \question{Soit $n\in\mathbb{N}$; montrer que pour tout entier $k\geq 1$ on a
$$n^k+kn^{k-1}\leq (n+1)^k.$$}
    \item \question{Soit $b$ un réel positif ou nul. Montrer par récurrence, que
pour tout $n\geq 1$ on a
$$ (1+b)^n\leq 1+\frac{nb}{1!}+\frac{(nb)^2}{2!}+...+\frac{(nb)^n}{n!}.$$}
\end{enumerate}
}
